%%
% BIThesis 本科毕业设计论文模板 —— 使用 XeLaTeX 编译 The BIThesis Template for Undergraduate Thesis
% This file has no copyright assigned and is placed in the Public Domain.
%%

\begin{conclusion}

  回归测试是保障软件质量的关键环节，其核心目标在于验证代码变更后既有功能的正确性不受影响。然而，传统的回归测试用例生成方式往往依赖人工编写或通用自动化工具，不仅耗时耗力，而且难以针对具体的代码变更进行精准覆盖。

  本文针对现有方法在变更感知能力、多智能体协作以及测试质量控制等方面的不足，提出并实现了一个面向代码变更的多智能体回归测试用例生成系统。该系统以LangGraph为编排框架，通过协调12个专业化的Agent节点，实现了从代码变更解析、影响范围推理、测试策略规划到测试用例生成、执行与评估的全链路自动化。

  在typer项目的typer.utils和typer.core两个模块上的实验结果表明，本文提出的多智能体系统在变更覆盖率、测试压缩率和执行效率等指标上均显著优于基线方法。特别是在变更感知的精确性方面，多智能体系统实现了93\%至96\%的测试压缩率，同时保持了78\%至96\%的执行时间缩减，且能够有效检测出Single-Agent方法遗漏的缺陷。消融实验验证了各Agent组件的不可替代性，其中LLM模块的移除会导致覆盖率下降56.8\%，TestRepair机制的缺失会使测试通过率从72.63\%降至37.86\%。

  综上所述，本文通过构建多智能体协作框架，将复杂的回归测试用例生成任务分解为多个专业化子任务，利用大规模语言模型的强大能力进行智能决策与生成，并在实验中取得了良好的效果。

  本文的主要贡献概括为以下几个方面：

  \textbf{（1）提出了一种基于多智能体协作的回归测试用例生成框架。} 不同于传统的单Agent或流水线式的测试生成方法，本文将测试生成过程建模为一个由12个专业化Agent节点构成的协作网络。每个Agent承担独立的职责，如代码变更解析、AST分析、影响范围推理、缺陷模式匹配等，通过LangGraph框架进行状态管理与任务调度。

  \textbf{（2）设计并实现了变更感知的测试用例生成机制。} 针对传统回归测试生成方法缺乏对代码变更精准感知的问题，本文提出了变更驱动的测试策略规划与用例生成方法。系统首先通过代码变更解析Agent识别变更的类型、范围与语义，然后由影响范围推理Agent定位受影响的函数与代码路径，最后由TestGenAgent针对性地生成覆盖这些变更的测试用例。

  \textbf{（3）构建了一套结合规则约束与LLM生成的混合测试策略。} 在测试用例生成过程中，本文设计了一套规则约束引擎与LLM协同工作的机制。规则引擎负责生成结构化的基础测试用例，确保覆盖基本的分支与语句；而LLM则在此基础上进行语义增强与边界条件补充。

  \textbf{（4）提出了基于TestRepairAgent的测试修复机制。} 本文引入TestRepairAgent对执行失败的测试用例进行自动修复，消融实验表明该机制可使测试通过率从37.86\%提升至72.63\%，增幅达47.9\%。

  \textbf{（5）通过全面的实验验证了系统的有效性与各组件的必要性。} 实验结果从多个维度证明了多智能体系统在覆盖率、效率与缺陷检测能力等方面的优势。

  尽管本文提出的多智能体回归测试用例生成系统取得了较好的实验效果，但仍存在一些局限性：

  \textbf{实验数据覆盖范围有限。} 本文的所有实验均基于typer\allowbreak 项目的\texttt{typer.utils}和\allowbreak\texttt{typer.core}两个模块展开，样本数量相对较少。

  \textbf{复杂模块上的覆盖率表现仍有提升空间。} 在\texttt{typer.core}模块上，分支覆盖率（79.51\%）和语句覆盖率（68.08\%）均低于Single-Agent方法（90.09\%，75.35\%）。

  \textbf{测试用例的生成通过率尚未达到100\%。} 尽管引入了TestRepairAgent，系统在typer.core模块上的测试通过率仅为70.1\%。

  \textbf{LLM生成的不确定性问题。} LLM的输出具有一定的随机性与不可控性。

  \textbf{生成时间在复杂场景下的优化需求。} 在涉及大规模代码变更或复杂逻辑的场景下，系统的端到端生成时间可能达到数十秒甚至更高。

  基于本文的研究工作与当前软件工程领域的发展趋势，未来可在以下几个方向进行深入探索与改进。

  \textbf{（1）系统与CI/CD流水线的深度集成。} 将测试用例生成系统接入持续集成/持续部署流程，在代码提交阶段自动触发测试生成，实现"提测前自动生成单元测试"的目标。

  \textbf{（2）扩展对更多编程语言与框架的支持。} 将系统的核心能力抽象为跨语言的测试生成框架，支持Java、Go、JavaScript、TypeScript等语言。

  \textbf{（3）优化LLM调用的效率与成本。} 引入模型缓存与复用机制，探索更轻量的模型或微调策略，设计智能调度策略。

  \textbf{（4）进一步提升测试通过率与用例质量。} 增强对测试依赖关系的自动分析与mock配置能力，引入更精确的断言生成与验证机制。

  \textbf{（5）实现增量测试执行与智能用例管理。} 研究基于代码变更图的增量测试选择技术，建立测试用例的生命周期管理机制。

  \textbf{（6）与开发工具链的深度集成。} 将测试生成能力集成到主流IDE、代码编辑器与协作平台中。

  总之，回归测试用例的自动化生成是一个具有重要实践价值的研究课题。本文在多智能体协作框架、变更感知生成与测试质量保障等方面进行了有益的探索，为该领域的研究与应用提供了新的思路。

\end{conclusion}
